\chapter{Implementation Decisions and Development Process}

\section{Technology Selection and Rationale}

\textbf{React Native with Expo Framework}

React Native with Expo was selected for cross-platform development targeting both iOS and Android with a single codebase within the 12-week timeline.

Key benefits included hot reload for rapid development iterations and access to established ML libraries. ML model integration requirements necessitated transitioning from Expo Go to custom development builds.

The app architecture centers around tab-based navigation defined in \texttt{app/(tabs)/\_layout.tsx}, organizing features into six main sections: Home, BirDex, Archive, Gallery, Account, and Settings. The bird logging workflow is separated into dedicated screens under \texttt{app/log/}, including manual entry (\texttt{manual.tsx}), photo capture (\texttt{photo.tsx}), and AI-powered real-time detection (\texttt{objectIdentCamera.tsx}).

\textbf{Machine Learning Framework Selection}

Local ML processing was implemented for privacy considerations. TensorFlow Lite models and MLKit packages avoided cloud infrastructure and API costs within academic project constraints.

Existing models were adapted rather than training custom models due to time limitations and educational focus on mobile development. BirdNET models were converted from .h5 to .tflite format and integrated whoBIRD Kotlin implementation concepts through the \texttt{services/ultraSimpleBirdClassifier.ts} service.

\textbf{Database Architecture Approach}

Database design combines SQLite for local storage with Firebase Firestore for cloud synchronization. This hybrid approach was motivated by learning objectives to gain experience with both local and cloud database management.

The implementation consists of two main database services: \texttt{services/database.ts} manages bird observation records, while \texttt{services/databaseBirDex.ts} handles the 30,000+ species database. Cloud synchronization is managed through \texttt{services/sync\_layer.ts}, providing optional Firebase integration. The offline-first design emerged from practical considerations during development, as it simplified testing and debugging when working with ML models that process large amounts of data locally.

\section{Camera and ML Integration Challenges}

\textbf{Camera Component Implementation}

Camera functionality integration with ML processing used React Native Vision Camera package for camera control and MLKit packages for object detection and image classification.

Initial real-time frame processing was unsuccessful as ML models couldn't process camera frames directly. An asynchronous pipeline in \texttt{app/log/objectIdentCamera.tsx} captures photos to temporary files, processes them through MLKit, and cleans up temporary files. Supporting utilities like \texttt{services/uriUtils.ts} handle file manipulation and image cropping operations.

\textbf{Processing Pipeline Development}

The image processing workflow follows these steps: capture a photo at reduced quality for ML processing, detect objects using MLKit, crop each detected object, classify the cropped images, and save results that meet the confidence threshold.

Proper error handling and temporary file cleanup proved crucial for app stability, particularly during extended testing sessions.

\section{Key Implementation Challenges}

\textbf{Development Environment Complexity}

The most significant challenge was transitioning from Expo Go to custom development builds when ML model integration required native module access, adding considerable complexity and requiring new build tools and debugging techniques.

\textbf{ML Package Selection}

We initially struggled with outdated or poorly maintained React Native ML packages. After research and testing, the MLKit packages from Infinite Red provided reliable object detection and image classification capabilities with good documentation and community support.

To ensure consistent user experience across the app, a comprehensive design system was developed with themed components in the \texttt{components/} directory. Core components like \texttt{ThemedView.tsx}, \texttt{ThemedText.tsx}, and \texttt{ThemedPressable.tsx} provide consistent styling and behavior throughout the application.

\textbf{Audio and Image Processing Coordination}

The original attempt to run image and audio ML processing concurrently led to systematic race conditions and file stability errors. The solution required implementing a sequential processing pipeline that handles one ML operation at a time, eliminating resource conflicts but requiring careful state management.

\textbf{File System Management}

Handling temporary files for ML processing required learning about mobile file system constraints and implementing proper cleanup mechanisms to prevent storage bloat during extended use.

State management across the application is handled through React Context providers in the \texttt{contexts/} directory. The \texttt{AuthContext.tsx} manages Firebase authentication with offline fallback, while \texttt{LogDraftContext.tsx} provides persistent draft storage for bird observations using AsyncStorage with auto-save functionality.




\section{Development Tools and Process}

\textbf{Development Environment}

Development used standard React Native tools including Metro bundler, TypeScript for type safety, and Expo CLI for build management. Custom development builds required EAS (Expo Application Services) for production builds.

Development workflow emphasized modular architecture with clear separation of concerns. Custom hooks like \texttt{hooks/useThemeColor.ts} provide centralized theme management, while \texttt{hooks/useBirdDexDatabase.ts} manages species database state across components.


\textbf{Development Assistance}

Development used online resources including documentation, community forums, and AI assistants for debugging and code generation, particularly when learning new frameworks or troubleshooting integration issues.

\section{Testing and Debugging Approach}

\textbf{Development Testing}

Testing used Android emulators and physical devices via USB debugging. Custom development builds eliminated Expo Go for rapid testing iterations.

Console logging throughout services tracked ML pipeline states and performance metrics, essential for debugging race condition issues that led to the unified pipeline solution. Development tools like \texttt{components/MemoryMonitor.tsx} helped track performance during testing sessions.

\textbf{Database Performance Optimization}

Managing the large BirDex dataset required strategic indexing, batch processing for translations, and efficient image storage formats for responsive performance with 33,000+ species records. The logging system required similar optimization for large media files and smooth cloud synchronization.

\textbf{Performance Considerations}

Performance testing focused on basic functionality verification, ensuring app responsiveness during ML processing and reasonable memory usage during extended use.

App performance varies with device capabilities, with older devices experiencing longer processing times for image and audio classification tasks.